\section{PRODUCTS OF SETS}

\begin{enumerate}
\def\labelenumi{\arabic{enumi}.}
\tightlist
\item
  Let E, F be two sets, which may or may not be distinct. The ordered pairs \((x, y)\) , whose first element x is any element of E and whose second element y is any element of F, are the elements of a new set, called the product of E by F, and denoted by \(E \times F\) ; E and F are called the factors of \(E \times F\) . Two ordered pairs are considered to be identical only if they have the same first element and the same second element; in other words, the relation `` \((x, y) = (x', y')\) '' is equivalent to the relation `` \(x = x'\) and \(y = y''\) ''. If z is any element of \(E \times F\) , the relation ``x is the first element of the ordered pair \(z''\) is a functional relation in x; it determines a mapping of \(E \times F\) onto E, which is called the first coordinate function, or the first projection, and is denoted by \(pr_{1}\) . Instead of saying ``x is the first element of the ordered pair \(z''\) , we also say ``x is the first coordinate (or projection) of \(z''\) or '' \(x = \mathrm{pr}_1z''\) . Similarly we define the second coordinate function, or second projection, which is a mapping of \(E \times F\) onto \(F\) , denoted by \(\mathrm{pr}_2\) .
\end{enumerate}

The relation `` \(x = \text{pr}_1 z\) and \(y = \text{pr}_2 z\) '' is equivalent to `` \(z = (x, y)\) ''.

The extension of the function \(pr_{1}\) to sets of subsets is denoted by the same symbol, in accordance with the general conventions, and is again called the first projection (here we do not use the term ``coordinate''); similarly for the extension of the second projection.

\begin{enumerate}
\def\labelenumi{\arabic{enumi}.}
\setcounter{enumi}{1}
\tightlist
\item
  A relation R between a generic element x of E and a generic element y of F is a property of the pair \((x, y)\) , and consequently defines a subset of the product \(E \times F\) , called the graph of R. Conversely, every subset A of \(E \times F\) is the graph of the relation \((x, y) \in A\) between x and y.
\end{enumerate}

Let A be a subset of E and let B be a subset of F. We denote by \(A \times B\) the subset of \(E \times F\) defined by the relation `` \(x \in A\) and \(y \in B\) '' between x and y.

\begin{enumerate}
\def\labelenumi{\arabic{enumi}.}
\setcounter{enumi}{2}
\tightlist
\item
  In the following propositions, X and X' denote arbitrary subsets of E, Y and Y' arbitrary subsets of F, and Z an arbitrary subset of E × F.
\end{enumerate}

\begin{enumerate}
\def\labelenumi{(\alph{enumi})}
\item
  The relation ``X × Y = φ'' is equivalent to ``X = φ or Y = φ''.
\item
  If \(X \times Y \neq \emptyset\) , the relation `` \(X \times Y \subset X' \times Y''\) '' is equivalent to `` \(X \subset X'\) and \(Y \subset Y''\) ''.
\item
  For all \(\mathbf{X},\mathbf{X}^{\prime},\mathbf{Y}\) we have
\end{enumerate}

\[
(\mathbf {X} \times \mathbf {Y}) \cup (\mathbf {X} ^ {\prime} \times \mathbf {Y}) = (\mathbf {X} \cup \mathbf {X} ^ {\prime}) \times \mathbf {Y}.\tag{22}
\]

\begin{enumerate}
\def\labelenumi{(\alph{enumi})}
\setcounter{enumi}{3}
\tightlist
\item
  For all \(\mathbf{X},\mathbf{X}^{\prime},\mathbf{Y},\mathbf{Y}^{\prime}\) we have
\end{enumerate}

\[
(\mathbf {X} \times \mathbf {Y}) \cap (\mathbf {X} ^ {\prime} \times \mathbf {Y} ^ {\prime}) = (\mathbf {X} \cap \mathbf {X} ^ {\prime}) \times (\mathbf {Y} \cap \mathbf {Y} ^ {\prime}).
\]

\begin{enumerate}
\def\labelenumi{(\alph{enumi})}
\setcounter{enumi}{4}
\tightlist
\item
  For all X, Y we have
\end{enumerate}

\[
\mathrm{pr} _ {1} ^ {- 1} (\mathbf {X}) = \mathbf {X} \times \mathbf {F}, \quad \mathrm{pr} _ {2} ^ {- 1} (\mathbf {Y}) = \mathbf {E} \times \mathbf {Y}.\tag{24}
\]

\begin{enumerate}
\def\labelenumi{(\alph{enumi})}
\setcounter{enumi}{5}
\tightlist
\item
  If \(\mathbf{Y} \neq \emptyset\) , then for all \(\mathbf{X}\) we have
\end{enumerate}

\[
\operatorname{pr} _ {1} (\mathbf {X} \times \mathbf {Y}) = \mathbf {X}.\tag{25}
\]

\begin{enumerate}
\def\labelenumi{(\alph{enumi})}
\setcounter{enumi}{6}
\tightlist
\item
  For all Z we have
\end{enumerate}

\[
Z \subset \operatorname{pr} _ {1} (Z) \times \operatorname{pr} _ {2} (Z).\tag{26}
\]

\begin{enumerate}
\def\labelenumi{(\alph{enumi})}
\setcounter{enumi}{7}
\tightlist
\item
  Let \(a\) be an element of \(\mathbf{E}\) . Then the mapping \((a, y) \to y\) of the set \(\{a\} \times \mathbf{F}\) onto \(\mathbf{F}\) (i.e., the restriction of \(\mathbf{pr}_2\) to the subset \(\{a\} \times \mathbf{F}\) ) is one-to-one.
\end{enumerate}

\begin{enumerate}
\def\labelenumi{\arabic{enumi}.}
\setcounter{enumi}{3}
\item
  The mapping

\[
(x, y) \rightarrow (y, x)\tag{27}
\]

is a one-to-one mapping of \(E \times F\) onto \(F \times E\) , and is called canonical. When E and F are the same set, the mapping (27) is called the canonical symmetry; it is then involutory. The elements \((x, y)\) of \(E \times E\) which are fixed under this symmetry are those which have the property x = y; the set \(\Delta\) of these elements is called the diagonal of \(E \times E\) . The mapping \(x \to (x, x)\) is a bijection of E onto \(\Delta\) , called the diagonal mapping of E into \(E \times E\) .

If Z is any subset of \(E \times F\) , the image of Z under the canonical mapping of \(E \times F\) onto \(F \times E\) is denoted by \(\overline{Z}^{1}\) . If X is any subset of E and Y any subset of F, then

\[
\overline{\mathbf {X} \times \mathbf {Y}}^{1} = \mathbf {Y} \times \mathbf {X}.
\]

If a relation R between x and y, considered as a property of the ordered pair \((x, y)\) , defines a subset A of \(E \times F\) , then the same relation, considered as a property of the pair \((y, x)\) , defines the subset \(\overline{A}^{1}\) of \(F \times E\) . R is equivalent to each of the relations \((x, y) \in A\) , \((y, x) \in \overline{A}^{1}\) . If E and F are the same set, the relation R and the corresponding subset A are said to be symmetric when \(A = \overline{A}^{1}\) . The diagonal \(\Delta\) (defined by the relation of equality) is symmetric. If Z is any subset of \(E \times E\) , then \(Z \cup \overline{Z}^{1}\) and \(Z \cap \overline{Z}^{1}\) are symmetric.

\end{enumerate}

\begin{enumerate}
\def\labelenumi{\arabic{enumi}.}
\setcounter{enumi}{4}
\tightlist
\item
  Let A be a subset of a set E, and let f be a mapping of A into a set F. The relation `` \(x \in A\) and \(y = f(x)\) '' between a generic element x of E and a generic element y of F defines a subset of \(E \times F\) called the graph of the function f.~If B is a subset of E which contains A, and if g is an extension (§ 2, no. 13) of f to B, then the graph of f is contained in the graph of g.
\end{enumerate}

Conversely, let C be a subset of \(\mathbf{E} \times \mathbf{F}\) such that, for each \(x \in \mathbf{E}\) , there exists at most one \(y \in \mathbf{F}\) such that \((x, y) \in \mathbf{C}\) . Then the relation \((x, y) \in \mathbf{C}\) between a generic element \(x\) of the set \(\operatorname{pr}_1(\mathbf{C})\) and a generic element \(y\) of \(\mathbf{F}\) is a functional relation in \(y\) and determines a mapping of \(\operatorname{pr}_1(\mathbf{C})\) into \(\mathbf{F}\) whose graph is \(\mathbf{C}\) .

The set of subsets C of \(E \times F\) which have the property ``for all \(x \in E\) there is at most one \(y \in F\) such that \((x, y) \in C\) '' (a set which is a subset of \(\mathfrak{P}(\mathbf{E} \times \mathbf{F})\) ) can therefore be put into one-to-one correspondence with the set of mappings of subsets of E into F.

Let \(f\) be an injective mapping of \(\mathbf{E}\) into \(\mathbf{F}\) , and let \(g\) be the inverse of \(f\) considered as a bijection of \(\mathbf{E}\) onto \(f(\mathbf{E})\) . If \(\mathbf{C}\) is the graph of \(f\) , then the graph of \(g\) is \(\overline{\mathbf{C}}^1\) .

\begin{enumerate}
\def\labelenumi{\arabic{enumi}.}
\setcounter{enumi}{5}
\tightlist
\item
  If \(f\) is a mapping of \(\mathbf{E}\) into \(\mathbf{F}\) , and \(\mathbf{C}\) is its graph in \(\mathbf{E} \times \mathbf{F}\) , the relation `` \(y = f(x)\) '' is equivalent to `` \((x, y) \in \mathbf{C}\) ''. The relation `` \(y \in f(\mathbf{X})\) '' is equivalent to ``there exists \(x\) such that \(x \in \mathbf{X}\) and \((x, y) \in \mathbf{C}\) ''.
\end{enumerate}

Now let K be any subset of E × F, and let X be any subset of E. Let K(X) denote the subset of F consisting of all elements y which satisfy the relation ``there exists x such that \(x \in X\) and \((x, y) \in K\) ''; this relation is thus equivalent to `` \(y \in K(X)\) ''. The mapping \(X \to K(X)\) of \(\mathfrak{P}(E)\) into \(\mathfrak{P}(F)\) is said to be defined by the subset K of E × F. Note that K(X) is the second projection of the set \(K \cap (X \times F)\) . If K is the graph of a mapping f of E into F, then the mapping \(X \to K(X)\) is the canonical extension of f to sets of subsets.

\begin{enumerate}
\def\labelenumi{\arabic{enumi}.}
\setcounter{enumi}{6}
\tightlist
\item
  If \(x\) is a generic element of \(\mathbf{E}\) , then \(x \to \mathrm{K}(\{x\})\) is a mapping of \(\mathbf{E}\) into \(\mathfrak{P}(\mathbf{F})\) whose value \(\mathrm{K}(\{x\})\) (denoted also by \(\mathrm{K}(x)\) , by abuse of language) is called the section of \(\mathbf{K}\) at \(x\) . The relation \((x, y) \in \mathbf{K}\) is equivalent to \(y \in \mathrm{K}(x)\) .
\end{enumerate}

Conversely, every mapping \(x \to \Phi(x)\) of \(E\) into \(\mathfrak{P}(F)\) can be obtained in this way; for the relation \(y \in \Phi(x)\) defines a subset \(K\) of \(E \times F\) , and \(\Phi(x)\) is precisely the section of \(K\) at \(x\) . The set \(\mathfrak{P}(E \times F)\) and the set of mappings of \(E\) into \(\mathfrak{P}(F)\) are thus in one-to-one correspondence.

\begin{enumerate}
\def\labelenumi{\arabic{enumi}.}
\setcounter{enumi}{7}
\tightlist
\item
  Every mapping \(\mathbf{X} \to \mathbf{K}(\mathbf{X})\) defined by a subset \(\mathbf{K}\) of \(\mathbf{E} \times \mathbf{F}\) has the following properties, which generalize those of the canonical extension of a mapping of \(\mathbf{E}\) into \(\mathbf{F}\) (§2, nos. 4 and 5):
\end{enumerate}

\begin{enumerate}
\def\labelenumi{(\alph{enumi})}
\item
  \(\mathbf{K}(\phi)=\phi.\)
\item
  ``X ⊂ Y'' implies ``K(X) ⊂ K(Y)''.
\item
  For all X, Y we have
\end{enumerate}

\begin{enumerate}
\def\labelenumi{(\arabic{enumi})}
\setcounter{enumi}{27}
\tightlist
\item
\end{enumerate}

\[
\mathbf {K} (\mathbf {X} \cup \mathbf {Y}) = \mathbf {K} (\mathbf {X}) \cup \mathbf {K} (\mathbf {Y}),\tag{29}
\]

\[
\mathbf {K} (\mathbf {X} \cap \mathbf {Y}) \subset \mathbf {K} (\mathbf {X}) \cap \mathbf {K} (\mathbf {Y}).
\]

If K and \(K'\) are two subsets of \(E \times F\) such that \(K \subset K'\) , then we have \(\mathrm{K}(\mathrm{X}) \subset \mathrm{K}'(\mathrm{X})\) for all \(X \subset E\) ; in particular, \(\mathrm{K}(x) \subset \mathrm{K}'(x)\) for all \(x \in E\) . Conversely, if \(\mathrm{K}(x) \subset \mathrm{K}'(x)\) for all \(x \in E\) , then \(K \subset K'\) .

\begin{enumerate}
\def\labelenumi{\arabic{enumi}.}
\setcounter{enumi}{8}
\tightlist
\item
  A relation between a generic element of E and a generic element of F defines a subset K of \(E \times F\) and a subset \(\overline{K}\) of \(F \times E\) , and hence a mapping \(X \to K(X)\) of \(\mathfrak{P}(E)\) into \(\mathfrak{P}(F)\) and a mapping \(Y \to \overline{K}(Y)\) of \(\mathfrak{P}(F)\) into \(\mathfrak{P}(E)\) .
\end{enumerate}

If \(\mathbf{K}\) is the graph of a mapping \(f\) of \(\mathbf{E}\) into \(\mathbf{F}\) , the mapping \(\mathbf{Y} \to \overline{\mathbf{K}}(\mathbf{Y})\) is the inverse extension of \(f\) .

It should be noted that the relations (18) and (19) do not generalize to the mappings \(\mathbf{X} \to \mathbf{K}(\mathbf{X})\) and \(\mathbf{Y} \to \overline{\mathbf{K}}^{1}(\mathbf{Y})\) , when \(\mathbf{K}\) is an arbitrary subset of \(\mathbf{E} \times \mathbf{F}\) .

\begin{enumerate}
\def\labelenumi{\arabic{enumi}.}
\setcounter{enumi}{9}
\tightlist
\item
  Let E, F, G be three sets, which may or may not be distinct, let A be a subset of \(E \times F\) and let B be a subset of \(F \times G\) . Then the elements \((x, z)\) of \(E \times G\) which have the property ``there exists \(y \in F\) such that \((x, y) \in A\) and \((y, z) \in B\) '' form a subset of \(E \times G\) , called the composition of B and A, and denoted by \(B \circ A\) , or simply BA when there is no risk of confusion. Here again, the order of composition is essential. ¶ The mapping \(X \to BA(X)\) of \(\mathfrak{P}(E)\) into \(\mathfrak{P}(G)\) is the composition of \(Y \to B(Y)\) and \(X \to A(X)\) ; in other words, for all \(X \subset E\) we have
\end{enumerate}

\[
\mathbf {B A} (\mathbf {X}) = \mathbf {B} (\mathbf {A} (\mathbf {X})).\tag{30}
\]

Let H be another set, not necessarily distinct from E, F, G, and let C be a subset of \(G \times H\) . Then we have \(\mathbf{C} \circ (\mathbf{B} \circ \mathbf{A}) = (\mathbf{C} \circ \mathbf{B}) \circ \mathbf{A}\) ; this set is also denoted by \(C \circ B \circ A\) (or simply CBA) and is called the composition of C, B, A taken in this order.

Let \(f\) be a mapping of \(\mathbf{E}\) into \(\mathbf{F}\) , and let \(g\) be a mapping of \(\mathbf{F}\) into \(\mathbf{G}\) . If \(\mathbf{A}\) (resp. \(\mathbf{B}\) ) is the graph of \(f\) (resp. \(g\) ), then the composition BA is the graph of the composite mapping \(g \circ f\) .

\begin{enumerate}
\def\labelenumi{\arabic{enumi}.}
\setcounter{enumi}{10}
\tightlist
\item
  We have
\end{enumerate}

\[
\overline{\mathbf {B} \circ \mathbf {A}}^{1} = \overline {\mathbf {A}}^{1} \circ \overline {\mathbf {B}}^{1}.\tag{31}
\]

Let A, A' be two subsets of E × F, and let B, B' be two subsets of F × G. Then the relation

\[
" \mathbf {A} \subset \mathbf {A} ^ {\prime} \text {   and   } \mathbf {B} \subset \mathbf {B} ^ {\prime}" \quad \text { implies } \quad " \mathbf {B} \circ \mathbf {A} \subset \mathbf {B} ^ {\prime} \circ \mathbf {A} ^ {\prime}".
\]

Let A be a subset of \(E \times F\) , \(\Delta\) the diagonal of \(E \times E\) , and \(\Delta'\) the diagonal of \(F \times F\) . Then we have

\[
\mathbf {A} \circ \Delta = \Delta^ {\prime} \circ \mathbf {A} = \mathbf {A}.\tag{32}
\]

\begin{enumerate}
\def\labelenumi{\arabic{enumi}.}
\setcounter{enumi}{11}
\tightlist
\item
  Let E, F, G be three sets, which may or may not be distinct. Their product \(E \times F \times G\) is the set of ordered triples \((x, y, z)\) , where \(x \in E\) , \(y \in F\) , and \(z \in G\) , the relation `` \((x, y, z) = (x', y', z')\) '' being equivalent to `` \(x = x'\) and \(y = y'\) and \(z = z''\) ''. The three mappings \((x, y, z) \to x\) , \((x, y, z) \to y\) , \((x, y, z) \to z\) of \(E \times F \times G\) onto E, F, G respectively are called the first, second, and third coordinate functions (or projections); similarly, for example, the projection with indices 1, 2 is the mapping
\end{enumerate}

\[
(x, y, z) \rightarrow (x, y)
\]

of \(\mathbf{E} \times \mathbf{F} \times \mathbf{G}\) onto \(\mathbf{E} \times \mathbf{F}\) , and is denoted by \(\mathbf{pr}_{1,2}\) .

The definitions and propositions of nos. 2, 3, and 4 generalize easily to the product of three sets.

Furthermore, instead of considering the product \(E \times F \times G\) of three sets, we may equivalently consider the product \((\mathbf{E} \times \mathbf{F}) \times \mathbf{G}\) , obtained by a double application of the operation of forming the product of two sets. In fact, \((x, y, z) \to ((x, y), z)\) is a one-to-one mapping of \(E \times F \times G\) onto \((\mathbf{E} \times \mathbf{F}) \times \mathbf{G}\) , called canonical. Similarly we define one-to-one mappings of \(E \times F \times G\) onto the set \(E \times (F \times G)\) , and onto all the sets obtained from \(E \times F \times G\) , \((\mathbf{E} \times \mathbf{F}) \times \mathbf{G}\) , and \(E \times (F \times G)\) by permuting the three letters E, F, G.

There are analogous definitions and properties for the product of more than three sets.

\begin{enumerate}
\def\labelenumi{\arabic{enumi}.}
\setcounter{enumi}{12}
\tightlist
\item
  If a function \(f\) , which takes its values in any set \(\mathbf{E}'\) , is defined on a product of three sets \(\mathbf{E}\) , \(\mathbf{F}\) , \(\mathbf{G}\) , it is said to be a function of three variables, each of which runs through one of the sets \(\mathbf{E}\) , \(\mathbf{F}\) , \(\mathbf{G}\) . The value of \(f\) at the element \((x, y, z)\) of \(\mathbf{E} \times \mathbf{F} \times \mathbf{G}\) is denoted by \(f(x, y, z)\) .
\end{enumerate}

Let \(a\) be any element of \(\mathbf{E}\) . Then \((y, z) \to f(a, y, z)\) is a mapping of \(\mathbf{F} \times \mathbf{G}\) into \(\mathbf{E}'\) , called a partial mapping (or function) determined by \(f\) corresponding to the value \(a\) of \(x\) ; it is also the composition of \(f\) and the mapping \((y, z) \to (a, y, z)\) of \(\mathbf{F} \times \mathbf{G}\) into \(\mathbf{E} \times \mathbf{F} \times \mathbf{G}\) .

Likewise, if \(b\) is an element of \(\mathbf{F}\) , then \(z \to f(a, b, z)\) is a mapping of \(\mathbf{G}\) into \(\mathbf{E}'\) , called the partial mapping determined by \(f\) corresponding to the values \(a, b\) of \(x, y\) .

Inversely, let g be a mapping of E into \(E'\) . Then \((x, y, z) \to g(x)\) is a mapping h of \(E \times F \times G\) into \(E'\) such that every partial mapping of E into \(E'\) determined by h, corresponding to any values of y and z, is identical with g. This fact is often expressed by saying that a function of an argument x can always be envisaged as a function of all the arguments which need to be considered at a given moment and which will, of course, include x.

\begin{enumerate}
\def\labelenumi{\arabic{enumi}.}
\setcounter{enumi}{13}
\tightlist
\item
  Let \(f, g, h\) be three mappings of \(\mathbf{E}\) into \(\mathbf{E}'\) , \(\mathbf{F}\) into \(\mathbf{F}'\) , \(\mathbf{G}\) into \(\mathbf{G}'\) , respectively. The mapping \((x, y, z) \to (f(x), g(y), h(z))\) of \(\mathbf{E} \times \mathbf{F} \times \mathbf{G}\) into \(\mathbf{E}' \times \mathbf{F}' \times \mathbf{G}'\) is denoted by \((f, g, h)\) or \(f \times g \times h\) and is called the extension of \(f, g, h\) to products. If all three of \(f, g, h\) are injective (resp. surjective, bijective), then \(f \times g \times h\) is injective (resp. surjective, bijective).
\end{enumerate}

In this and the previous subsection we have considered only the case of three sets merely to fix the ideas; analogous considerations hold for any finite number of sets.

